% =============================================================================
% ANEXO X — Catálogo de Códigos de Error y Respuesta
% =============================================================================
\chapter{Catálogo de Códigos de Error y Respuesta}
\label{cap:anexo_errores}

Códigos de estado HTTP empleados por la API y su semántica en el envoltorio
\comando{ApiResponse} (Sección~\ref{sec:anexo_ep_envelope}), junto con el mapeo de errores de IA.\par

\section{Códigos HTTP}
\label{sec:anexo_err_http}
\renewcommand{\arraystretch}{1.3}
\fontsize{9}{10.8}\selectfont\sffamily
\noindent
\begin{tabularx}{\textwidth}{|l|X|}
\hline
\rowcolor{colorPrimario}
\textcolor{white}{\textbf{Código}} & \textcolor{white}{\textbf{Semántica}} \\ \hline
200 OK & Operación exitosa con \textit{payload} en \comando{data}. \\ \hline
\rowcolor{grisTecnico}
201 Created & Recurso creado (registro, proyecto, skill). \\ \hline
400 Bad Request & Error de validación; \comando{errors} con los mensajes. \\ \hline
\rowcolor{grisTecnico}
401 Unauthorized & Token ausente, inválido o expirado. \\ \hline
403 Forbidden & Rol insuficiente (RBAC / \comando{RoleScopeGuard}). \\ \hline
\rowcolor{grisTecnico}
404 Not Found & Recurso o RPC inexistente. \\ \hline
409 Conflict & Conflicto (p. ej. email ya registrado). \\ \hline
\rowcolor{grisTecnico}
429 Too Many Requests & Límite de tasa superado (\textit{rate limiting}). \\ \hline
503 Service Unavailable & Indisponibilidad \textit{upstream} (IA). \\ \hline
\end{tabularx}\normalfont\normalsize
\par\vspace{0.3cm}

\section{Mapeo de Errores de IA (\texttt{AiServiceException})}
\label{sec:anexo_err_ia}
El manejador centraliza los fallos \textit{upstream} de Gemini para no exponer 500 genéricos,
según la tabla~\ref{tab:anexo_err_ia}.

\begin{table}[!ht]
    \centering
    \renewcommand{\arraystretch}{1.3}
    \fontsize{9}{10.8}\selectfont\sffamily
    \begin{tabularx}{\textwidth}{|l|l|X|}
        \hline
        \rowcolor{colorPrimario}
        \textcolor{white}{\textbf{Upstream}} & \textcolor{white}{\textbf{Devuelto}} & \textcolor{white}{\textbf{Interpretación}} \\ \hline
        429 & 429 & Saturación del proveedor; reintento diferido. \\ \hline
        \rowcolor{grisTecnico}
        5xx & 503 & Indisponibilidad temporal del servicio de IA. \\ \hline
        otros & 502 & Error de pasarela no clasificado. \\ \hline
    \end{tabularx}
    \caption{Mapeo de errores \textit{upstream} de IA (\comando{AiServiceException}).}
    \label{tab:anexo_err_ia}
\end{table}

\clearpage
\begin{NotaTecnica}
El manejador global de excepciones de \comando{shared/} traduce las excepciones de validación y
de negocio a respuestas \comando{ApiResponse} consistentes, garantizando que el cliente reciba
siempre la misma estructura, con \comando{success=false} y la lista \comando{errors} poblada.
\end{NotaTecnica}

\clearpage
